% !TeX encoding = UTF-8
\documentclass[variant=apuntes,cover=institutional,mode=screen]{ua-engineering-v6-article}
% !TeX encoding = UTF-8
\UASetUniversity{Universidad Austral}
\UASetFaculty{Facultad de Ingeniería}
\UASetDepartment{Departamento de Computación}
\UASetCampus{Pilar}
\UASetCareer{Ingeniería Informática}
\UASetCommission{Comisión A}
\UASetProfessor{Equipo docente}
\UASetSemester{Segundo cuatrimestre 2026}
\UASetCourse{Sistemas Distribuidos}
\UASetDate{2026}


\UASetClassNumber{13}
\UASetClassTitle{Confiabilidad operativa}
\UASetDocKind{Apuntes}

\title{Clase 13 --- Confiabilidad operativa: SLI, SLO, SLA, error budgets y operación guiada por objetivos}
\author{\UAProfessor}
\date{\UADate}

\begin{document}
\maketitle

\section{Introducción}

Una vez que el sistema ya puede observarse, resistir fallas y someterse a experimentación controlada, aparece una pregunta inevitable: ¿cómo decidimos si está funcionando lo suficientemente bien?

Responder “anda bien” o “anda mal” sin criterios explícitos conduce a operación reactiva, discusiones subjetivas y prioridades inconsistentes. Por eso, la confiabilidad operativa moderna se apoya en objetivos formales.

\begin{quote}
La confiabilidad no se define por perfección absoluta, sino por el cumplimiento de objetivos explícitos relevantes para el usuario.
\end{quote}

\section{Service Level Indicators (SLI)}

\subsection{Definición}

Un Service Level Indicator (SLI) es una métrica que captura una propiedad relevante del servicio desde la perspectiva del usuario o del comportamiento observable del sistema.

\subsection{Ejemplos}

\begin{itemize}
\item porcentaje de requests exitosas;
\item latencia p95 o p99;
\item porcentaje de jobs completados;
\item retraso máximo tolerado en un pipeline de eventos;
\item porcentaje de lecturas correctas o dentro de tiempo esperado.
\end{itemize}

\subsection{Qué hace bueno a un SLI}

Un buen SLI debe ser:

\begin{itemize}
\item \textbf{relevante}: reflejar algo que realmente importe a usuarios o al producto;
\item \textbf{medible}: poder calcularse con datos confiables;
\item \textbf{estable}: no depender de señales erráticas o mal definidas;
\item \textbf{accionable}: permitir tomar decisiones operativas o de ingeniería.
\end{itemize}

\subsection{Ejemplo de mal SLI}

Un indicador como “CPU promedio del cluster” puede ser útil para operación, pero no necesariamente es un buen SLI si no expresa una propiedad observable del servicio para el usuario.

\section{Service Level Objectives (SLO)}

\subsection{Definición}

Un Service Level Objective (SLO) es un objetivo cuantitativo sobre uno o más SLI.

Ejemplos:
\begin{itemize}
\item 99.9\% de requests exitosas en 30 días;
\item p95 de latencia menor a 300 ms;
\item 99.95\% de jobs procesados en menos de 10 segundos.
\end{itemize}

\subsection{Propósito}

El SLO convierte una expectativa vaga en un criterio operativo explícito. Esto permite:
\begin{itemize}
\item evaluar si el servicio está cumpliendo;
\item priorizar trabajo de confiabilidad;
\item decidir si conviene desplegar, frenar o invertir en mejoras;
\item comunicar con claridad el nivel de servicio esperado.
\end{itemize}

\subsection{Ventana temporal}

Todo SLO necesita una ventana, por ejemplo:
\begin{itemize}
\item 7 días;
\item 30 días;
\item trimestre;
\item rolling window.
\end{itemize}

La ventana importa porque cambia la interpretación del incumplimiento: una mala hora no significa lo mismo que una degradación sostenida durante semanas.

\section{Service Level Agreements (SLA)}

\subsection{Definición}

Un Service Level Agreement (SLA) es un compromiso formal y externo, a menudo contractual, asociado a uno o más objetivos de nivel de servicio.

\subsection{Relación con SLI y SLO}

\begin{itemize}
\item \textbf{SLI}: qué se mide;
\item \textbf{SLO}: qué objetivo se fija;
\item \textbf{SLA}: qué se promete formalmente a terceros.
\end{itemize}

\subsection{Importancia práctica}

Un equipo puede operar con SLOs internos más exigentes que el SLA externo, para reservar margen de seguridad.

\section{Error Budgets}

\subsection{Definición}

Si el SLO no exige 100\% de cumplimiento, entonces existe una fracción de falla tolerada. Esa fracción es el \emph{error budget}.

Ejemplo:
\begin{itemize}
\item SLO = 99.9\% de disponibilidad mensual;
\item error budget = 0.1\% del tiempo o requests del período.
\end{itemize}

\subsection{Interpretación}

El error budget no es permiso para degradar arbitrariamente el sistema, sino una forma de balancear:
\begin{itemize}
\item confiabilidad;
\item ritmo de cambio;
\item riesgo aceptable;
\item velocidad de innovación.
\end{itemize}

\subsection{Uso operativo}

Si el sistema consume demasiado rápido su error budget, el equipo puede decidir:
\begin{itemize}
\item frenar despliegues;
\item priorizar deuda técnica;
\item investigar degradación;
\item endurecer validaciones;
\item reconfigurar alertas o mecanismos de resiliencia.
\end{itemize}

\section{Operación guiada por objetivos}

\subsection{Sin objetivos explícitos}

Operar sin SLOs conduce a:
\begin{itemize}
\item incidentes juzgados subjetivamente;
\item discusiones ambiguas entre producto e ingeniería;
\item dificultad para priorizar mantenimiento vs nuevas features;
\item alertas ruidosas o inconsistentes.
\end{itemize}

\subsection{Con objetivos explícitos}

Operar con SLOs permite:
\begin{itemize}
\item tener un lenguaje compartido;
\item convertir confiabilidad en una variable administrable;
\item conectar decisiones técnicas con impacto de producto;
\item establecer umbrales razonables de degradación.
\end{itemize}

\section{Diseño de buenos SLOs}

\subsection{Principios}

Un buen SLO debe ser:
\begin{itemize}
\item comprensible;
\item medible;
\item centrado en la experiencia del usuario;
\item suficientemente exigente para importar;
\item suficientemente realista para ser sostenible.
\end{itemize}

\subsection{Errores frecuentes}

\begin{itemize}
\item definir SLOs imposibles de medir;
\item usar promedios cuando los percentiles son lo relevante;
\item medir infraestructura en vez de servicio;
\item mezclar muchas propiedades en un único número opaco;
\item fijar objetivos irrealistas que fuerzan sobreingeniería.
\end{itemize}

\section{Ejemplos}

\subsection{API pública}
\begin{itemize}
\item SLI: porcentaje de requests exitosas;
\item SLO: 99.95\% en rolling window de 30 días.
\end{itemize}

\subsection{Servicio de checkout}
\begin{itemize}
\item SLI: porcentaje de checkouts completados exitosamente;
\item SLO: 99.9\% mensual;
\item error budget: 0.1\%.
\end{itemize}

\subsection{Pipeline de eventos}
\begin{itemize}
\item SLI: porcentaje de eventos procesados en menos de 5 segundos;
\item SLO: 99.5\% semanal.
\end{itemize}

\section{Tensiones reales}

Elegir SLOs obliga a enfrentar preguntas difíciles:
\begin{itemize}
\item ¿qué tan rápido debe responder el sistema para que siga siendo útil?;
\item ¿cuánta degradación puede tolerarse?;
\item ¿qué servicio merece más presupuesto de confiabilidad?;
\item ¿qué tipo de error es más costoso para el negocio?
\end{itemize}

Esto conecta la confiabilidad con decisiones de arquitectura, resiliencia y producto.

\section{Modelo del curso}

Podemos expresar este problema como:

\[
D = (P, F, G, S, T)
\]

Ejemplo:
\begin{itemize}
\item $P$: operar un servicio con criterio explícito de calidad;
\item $F$: fallas parciales, degradaciones, incidentes, ruido observacional;
\item $G$: cumplimiento de nivel de servicio acordado;
\item $S$: SLI + SLO + error budgets + políticas operativas;
\item $T$: disciplina organizacional, costo de medición, decisiones más visibles y trazables.
\end{itemize}

\section{Conclusión}

La confiabilidad operativa madura no se basa en intuiciones del tipo “esto parece estable”. Se basa en métricas relevantes, objetivos explícitos y presupuestos de error que permitan tomar decisiones consistentes.

\begin{uakeybox}
Confiabilidad operativa significa administrar explícitamente cuánto falla el sistema, cuánto puede tolerarse y cuándo el costo del error obliga a cambiar el ritmo de cambio.
\end{uakeybox}


\section{Síntesis razonada de la clase}
Esta unidad debe leerse como parte de una cadena de decisiones de diseño y no como un catálogo de mecanismos. Para estudiar el tema con profundidad conviene reconstruir cuatro preguntas: \emph{qué problema aparece}, \emph{qué garantía se desea}, \emph{qué mecanismo la aproxima} y \emph{qué costo introduce}. En cada ejemplo debe identificarse además el modelo de falla implícito.

\begin{uakeybox}
Una respuesta técnicamente madura no termina al nombrar un algoritmo o patrón: declara sus supuestos, la propiedad que preserva y el escenario en el que deja de ser suficiente.
\end{uakeybox}

\section{Bibliografía guiada y ruta de profundización}
Para profundizar esta clase se recomienda: Google SRE Book y SRE Workbook, capítulos de SLOs y error budgets.

La lectura debe hacerse con preguntas concretas: (1) ¿qué modelo de sistema asume el autor?, (2) ¿qué propiedad demuestra o persigue?, (3) ¿qué falla o anomalía motiva el mecanismo?, y (4) ¿qué trade-off queda abierto?

\section{Conexión con el Trabajo Final Integrador}
El grupo debe señalar qué parte de su sistema utiliza los conceptos de esta unidad, justificar por qué son necesarios y documentar una alternativa descartada. Cuando el contenido todavía no corresponda a la etapa vigente, debe registrarse como hipótesis o decisión pendiente.

\section{Autoevaluación conceptual}
\begin{enumerate}
\item ¿Cuál es el problema distribuido concreto que motiva esta clase?
\item ¿Qué propiedad de seguridad y qué propiedad de progreso están en juego?
\item ¿Qué supuesto de red, tiempo o falla resulta decisivo?
\item ¿Qué trade-off aceptaría en un sistema real y cuál no aceptaría en un dominio crítico?
\item ¿Cómo observaría experimentalmente que la solución se comporta como espera?
\end{enumerate}

\section{Profundización autocontenida v9}
\subsection{Problema, límite e invariante}
El núcleo de esta clase es SLI, SLO, SLA, error budgets, burn rates y operación guiada por objetivos. Para estudiarlo de manera autónoma conviene separar tres planos. Primero, el \emph{problema observable}: qué comportamiento incorrecto puede ver un cliente o un operador. Segundo, el \emph{límite del modelo}: qué información, sincronía o coordinación no está disponible. Tercero, el \emph{invariante}: qué propiedad no debe violarse aunque el sistema pierda progreso temporalmente.

El modelo de fallas relevante incluye SLI mal elegido, ventana incorrecta, budget consumido y alerta ruidosa. Una solución debe describirse como protocolo y no solo como nombre de tecnología: actores, estados, mensajes, condición de éxito, condición de aborto y estado visible después de una interrupción. La evidencia esperada es SLI calculado, SLO, budget, burn rate y política de decisión.

\subsection{Método de análisis}
Ante un caso nuevo, siga esta secuencia:
\begin{enumerate}
\item Identifique actores, dato o recurso compartido y frontera de confianza.
\item Escriba una ejecución nominal y una ejecución adversarial.
\item Distinga seguridad (lo malo no ocurre) de progreso (algo bueno finalmente ocurre).
\item Declare qué supuesto permite avanzar: mayoría, deadline, sincronía parcial, operación idempotente, almacenamiento durable u otro.
\item Compare una alternativa más simple y una más fuerte; registre qué métrica o experimento haría cambiar la decisión.
\end{enumerate}

\subsection{Caso transferible}
Considere una plataforma de reservas que recibe operaciones duplicadas, pierde conectividad entre regiones y debe sostener un camino crítico sin ocultar el estado real al usuario. Aplique el mecanismo de esta clase a una única decisión concreta. Una respuesta madura no intenta resolver todo: delimita la operación, formula la garantía y explica la degradación esperada cuando el supuesto central deja de cumplirse.

\subsection{Errores de razonamiento frecuentes}
\begin{itemize}
\item confundir una propiedad con el producto que suele implementarla;
\item describir el camino feliz sin recorrer una falla intermedia;
\item afirmar ``exactly-once'', ``alta disponibilidad'' o ``consistencia'' sin definir alcance observable;
\item medir solo throughput aunque la hipótesis trate sobre semántica o recuperación;
\item presentar la complejidad añadida como beneficio sin compararla con la alternativa más simple.
\end{itemize}

\section{Lectura guiada}
Google SRE y SRE Workbook, capítulos de SLOs y error budgets.

Durante la lectura, registre: modelo de sistema, propiedad perseguida, contraejemplo que motiva el mecanismo, costo principal y una condición bajo la cual no lo elegiría. Estas cinco anotaciones convierten la bibliografía en una herramienta de diseño y no en una lista para memorizar.

\section{Aplicación al Trabajo Final Integrador}
El equipo debe producir un ADR breve con la decisión asociada a esta clase. Debe contener: contexto, dos alternativas, supuesto decisivo, garantía elegida, trade-off, evidencia pendiente y fecha de revisión. La evidencia mínima esperada para esta unidad es: SLI calculado, SLO, budget, burn rate y política de decisión.

\section{Autoevaluación avanzada}
\begin{enumerate}
\item ¿Puedo explicar la clase sin nombrar primero una tecnología?
\item ¿Puedo construir una ejecución que viole la garantía si el mecanismo se configura mal?
\item ¿Puedo distinguir seguridad, progreso y experiencia del usuario?
\item ¿Puedo proponer una medición que valide la propiedad, no solo el rendimiento?
\item ¿Puedo defender cuándo una solución más simple sería suficiente?
\end{enumerate}

\end{document}
